\section{奇异分解和伪逆}

逆矩阵的适用有严格限制：矩阵为方阵，矩阵满秩。伪逆则将“逆”拓展到了一般矩阵中。

\subsection{伪逆的定义}

\begin{BoxDefinition}[伪逆]
    设$\vb{A}\in\R^{m\times n}$，定义伪逆$\vb{A}^{\dagger}\in\R^{n\times m}$为（Pseudo Inverse）
    \begin{Equation}[伪逆]
        \vb{A}^{\dagger}=\vb{V}_1\tilde{\vb*{\Sigma}}\vb{U}_1^T
    \end{Equation}
\end{BoxDefinition}

伪逆满足以下四个条件，实际上，这也可以视为另外一种定义伪逆的方式。

\begin{BoxProperty}[伪逆的Moore-Penrose条件]
    伪逆$\vb{A}^{\dagger}$满足以下四个条件
    \begin{itemize}
        \item $\vb{A}\vb{A}^{\dagger}\vb{A}=\vb{A}$
        \item $\vb{A}^{\dagger}\vb{A}\vb{A}=\vb{A}$
        \item $\vb{A}\vb{A}^{\dagger}$是对称的
        \item $\vb{A}^{\dagger}\vb{A}$是对称的
    \end{itemize}
\end{BoxProperty}

伪逆一般来说$\vb{A}\vb{A}^{\dagger}\neq\vb{I}$和$\vb{A}^{\dagger}\vb{A}=\neq\vb{I}$，但在某些特殊情况下会成立。

\subsection{伪逆的性质}

\begin{BoxProperty}[伪逆的自反性]
    伪逆满足自反性
    \begin{Equation}[伪逆的自反性]
        (\vb{A}^{\dagger})^{\dagger}=\vb{A}
    \end{Equation}
\end{BoxProperty}

\begin{BoxProperty}[伪逆和转置]
    伪逆和转置可以交换顺序
    \begin{Equation}[伪逆和转置]
        (\vb{A}^T)^{\dagger}=(\vb{A}^{\dagger})^T
    \end{Equation}
\end{BoxProperty}

\begin{BoxProperty}[伪逆和数乘]
    伪逆中的数乘可以提出
    \begin{Equation}[伪逆和数乘]
        (a\vb{A})^{\dagger}=a^{-1}\vb{A}^{\dagger}
    \end{Equation}
\end{BoxProperty}

\begin{BoxProperty}[伪逆不改变秩]
    伪逆不改变秩
    \begin{Equation}[伪逆不改变秩]
        \rank(\vb{A})=\rank(\vb{A}^{\dagger})=\rank(\vb{A}^{\dagger}\vb{A})=\rank(\vb{A}\vb{A}^{\dagger})
    \end{Equation}
\end{BoxProperty}

关于伪逆和转置，这两者之间存在许多组的恒等式。

\begin{BoxProperty}[伪逆和转置的恒定式1]
    伪逆和转置满足以下恒等式
    \begin{Gather}[伪逆和转置的恒定式1.]
        (\vb{A}\vb{A}^T)^{\dagger}=(\vb{A}^T)^{\dagger}\vb{A}^{\dagger}\id{a}\\
        (\vb{A}^T\vb{A})^{\dagger}=\vb{A}^{\dagger}(\vb{A}^T)^{\dagger}\id{b}
    \end{Gather}
\end{BoxProperty}

\begin{BoxProperty}[伪逆和转置的恒定式2]
    伪逆和转置满足以下恒等式
    \begin{Gather}[伪逆和转置的恒定式2.]
        \vb{A}^{\dagger}=(\vb{A}^T\vb{A})^{\dagger}\vb{A}^T=\vb{A}^T(\vb{A}\vb{A}^T)^{\dagger}\id{a}
    \end{Gather}
\end{BoxProperty}

\begin{BoxProperty}[伪逆和转置的恒定式3]
    伪逆和转置满足以下恒等式
    \begin{Gather}[伪逆和转置的恒定式3.]
        \vb{A}\vb{A}^{\dagger}=\vb{A}\vb{A}^T(\vb{A}\vb{A}^T)^{\dagger}=(\vb{A}\vb{A}^T)^{\dagger}\vb{A}\vb{A}^T\id{a}\\
        \vb{A}^{\dagger}\vb{A}=\vb{A}^T\vb{A}(\vb{A}^T\vb{A})^{\dagger}=(\vb{A}^T\vb{A})^{\dagger}\vb{A}^T\vb{A}\id{b}
    \end{Gather}
\end{BoxProperty}

伪逆可以对任意的$\vb{A}\in\R^{m\times n}$定义，但如果$\vb{A}$能做到列满秩或行满秩，则有更简单的定义。

\begin{BoxProperty}[列满秩和左逆]
    若$\vb{A}$是列满秩的，则其具有以下性质，此时可以称为左逆（Left Inverse）
    \begin{Equation}[列满秩和左逆]
        \vb{A}^{\dagger}=(\vb{A}^T\vb{A})^{-1}\vb{A}^T\qquad
        \vb{A}^{\dagger}\vb{A}=\vb{I}
    \end{Equation}
\end{BoxProperty}

\begin{BoxProperty}[行满秩和右逆]
    若$\vb{A}$是行满秩的，则其具有以下性质，此时可以称为右逆（Right Inverse）
    \begin{Equation}[行满秩和右逆]
        \vb{A}^{\dagger}=\vb{A}^T(\vb{A}\vb{A}^T)^{-1}\qquad
        \vb{A}\vb{A}^{\dagger}=\vb{I}
    \end{Equation}
\end{BoxProperty}
